\documentclass{article}
\usepackage[UTF8]{ctex}
\usepackage{amsmath,amssymb,amsfonts}
\usepackage{enumitem}
\usepackage{geometry}
\geometry{a4paper, margin=1in}
% === 错题本排版优化（不参与题目标识） ===
\setlength{\parindent}{0pt}
\setlength{\parskip}{0.35em plus 0.10em minus 0.08em}
\setlength{\emergencystretch}{3em}
\setlist{itemsep=0.25em, topsep=0.40em, parsep=0pt}
\allowdisplaybreaks[2]
% === 错题本排版优化结束 ===
\begin{document}
\section*{中南大学2025年数学分析真题}
\section*{145．中南大学2025年数学分析试卷}

\subsection{题目1}
一．计算题．每题10分，共40分．

\begin{itemize}
  \item[（1）] \(\lim _{x \rightarrow 0^{+}} \frac{\sin (\tan x)-\tan (\sin x)}{\sin ^{4} x}\) ．
  \item[（2）] \(\int_{0}^{2} \frac{\mathrm{~d} x}{e^{x}+e^{-x}}\) ．
  \item[（3）] \(\lim _{n \rightarrow \infty}\left(1-\frac{1}{n}+\frac{1}{n^{2}}\right)^{1-n}\) ．
  \item[（4）] \(\int_{0}^{1} \frac{\cos x \cos (\sin x)}{\sin (\sin x)+\cos (\sin x)} \mathrm{d} x\) ．
\end{itemize}

\subsection{题目2}
二．（15分）确定如下级数 \(\sum_{n=0}^{\infty} \frac{2 n+1}{3^{n+2}} x^{2 n}\) 的收敛域，并计算在收敛域内的和函数．

\subsection{题目3}
三．（15 分）证明：函数 \(y=\frac{x}{1+x^{4}}\) 在实数域 \(\mathbb{R}\) 上有界且一致连续．

\subsection{题目4}
四．（15 分）计算曲线 \(x=e^{t} \cos t, y=e^{t} \sin t\) 在 \(t \in[0, \pi]\) 的长度．

\subsection{题目5}
五．（15 分）求由方程 \(x^{2}+2 x y+2 y^{2}=1\) 所确定的隐函数的极值．

\subsection{题目6}
六。（15 分）证明：若函数 \(u=f(x, y)\) 满足拉普拉斯方程 \(u_{x x}+u_{y y}=0\) ，则函数 \(v=\) \(f(x-\sqrt{3} y, \sqrt{3} x+y)\) 也满足拉普拉斯方程．

\subsection{题目7}
七．（15 分）设 \(f:(0,1) \rightarrow \mathbb{R}\) 满足以下条件（＊）：存在 \(L>0, \alpha \in(0,1)\) ，使得对任意的 \(x, y \in(0,1)\) ，始终有

\[
|f(x)-f(y)| \leq L|x-y|^{\alpha} .
\]

（1）（6 分）证明：\(f\) 在 \((0,1)\) 上连续．\\
（2）（6 分）证明：\(f\) 在 \((0,1)\) 上一致连续．\\
（3）（3 分）如果函数 \(f\) 满足条件（＊），则称其为 \(C^{\alpha}\) 函数，请举出一个例子，满足对于 \(0<\beta<\alpha\) ，该函数是 \(C^{\beta}\) 函数，但它不是 \(C^{\alpha}\) 函数．

\subsection{题目8}
八．（20 分）解答如下问题：\\
（1）（10 分）计算 \(f(x)=\frac{2 x \sin \theta}{1-2 x \cos \theta+x^{2}}\) 在 \(x=0\) 处的泰勒展开式。\\
（2）（10 分）设 \(|x|<1\) ，计算 \(\int_{0}^{\pi} \ln \left(1-2 x \cos \theta+x^{2}\right) \mathrm{d} \theta\) ．
\end{document}
